\documentclass[11pt,a4paper]{article}
\usepackage[spanish]{babel}
\usepackage[utf8]{inputenc}
\usepackage{graphicx}
\usepackage[left=2.5cm,top=3cm,right=2.5cm,bottom=3cm,bindingoffset=0.5cm]{geometry}
\usepackage{macros/AEDLogica, macros/AEDEspecificacion, macros/AEDTADs}
\usepackage{caratula}
\usepackage{amsmath}
\usepackage{amssymb}
\usepackage{amsfonts}

\titulo{Trabajo práctico 1: TAD AirMiles.}
\subtitulo{Especificación del TAD AirMiles: sistema de fidelización de una aerolinea.}

\fecha{\today}

\materia{Algoritmos y Estructuras de Datos}
\grupo{Los4Bugueados}

\integrante{Campos, Felipe}{356/20}{felippe.1114@gmail.com}
% \integrante{Apellido, Nombre2}{002/01}{email2@dominio.com} %
\integrante{Portaluppi, Alejandro}{51/19}{alejandro\_portaluppi@outlook.com}
\integrante{Chopite, Julio}{617/24/01}{email2@dominio.com}
\integrante{Rychter, Ariel}{1294/21}{arirychter@gmail.com}

% Declaramos donde van a estar las figuras
% No es obligatorio, pero suele ser comodo
\graphicspath{{../static/}}

% Asi pueden escribir nuevos comandos.
% Este por ejemplo asegura q los nombres
% que figuren con una tipografia diferenciada
\newcommand{\Tipo}[1]{\mathsf{#1}}
% la sintaxis es \newcommand{\nombreDeLaMacro}[cantidadDeParametros]{Lo que va ser remplazado por el macro}
\newcommand{\norm}[1]{\vert #1\vert}


\begin{document}

\maketitle

\section{Macros para escribir en LaTeX}

Para facilitar la escritura del trabajo práctico en LaTeX, les proveemos de varias macros que les permitirán escribir especificaciones en lógica de primer orden.
Por ejemplo, podrán hacer uso de comandos que definen los conectores lógicos: $\land, \lor, \implica$ y también los operadores de la lógica trivaluada $\yLuego$, $\oLuego$, $\implicaLuego$.

Además, disponen de comandos para cuantificadores: $\paraTodo{x}{\mathbb{Z}}{x \geq 0 \lor x < 0}$, $\existe{x}{\mathbb{Z}}{x \geq 0}$. A estos comandos se les puede agregar un \texttt{Largo} al final para que ocupen múltiples renglones.

\section{Ejemplos de uso de las macros}


\begin{tad}{EjemploDeTAD}


\obs{secuenciaDeCaracteres}{\seq{\cha}}
\\
\obs{conjuntoDeTuplas}{\conj{\tupla{nombre: \seq{\cha}, apellido: \seq{\cha}}}}
\\


\begin{proc}{crearTAD}
{
\In caracteres: \seq{\cha},
\In secuenciaDeNombres: \seq{\tupla{nombre: \seq{\cha}, apellido: \seq{\cha}}}
}{
\Tipo{EjemploDeTAD}
}
    \requiere{predicadoUnilinea(caracteres)}
    \asegura{res.secuenciaDeCaracteres = caracteres}
    \aseguraLargo{ \norm{res.conjuntoDeTuplas} = \norm{secuenciaDeNombres} \land predicadoMultilinea(secuenciaDeNombres, res.conjuntoDeTuplas) }
\end{proc}

\predLargo{predicadoMultilinea}{s: \seq{\cha}, conjunto: \conj{\tupla{\seq{\cha}, \seq{\cha}}}}{
    \paraTodo{i}{\Z}{0 \leq i < \norm{s}
    \implicaLuego s[i] \in c}
}

\pred{predicadoUnilínea}{s: \seq{\cha}}{
    \norm{s} > 0 \ (\forall i:\mathbb{Z})
}

\aux{unAuxiliar}{k: \Z,l: \Z}{\Z}{\norm{k-l}}

\end{tad}

\section{Especificacion del TAD AirMiles}

\noindent TipoDeTransaccion ES \{ Transferencia, Acumular, Canjear \}

\noindent Trans ES $\struct{origen: \Z, destino: \Z, cantidad: \Z}$

%\noindent Trans ES $\struct{origen: \Z, destino: \Z, cantidad: \Z, tipoDeTransaccion}$

\noindent Categorias ES \{ Bronce, Plata, Oro, Platino \}

\noindent Acumulacion ES $\struct{distanciaVolada:\R, promo:\tupla{nombrePromo, factor}, millasGanadas:\R}$

%ESTE ES EL NUEVO OBSERVADOR SUGERIDO - Julio
% \noindent Acumulacion ES $\struct{distanciaVolada:\R, promo:\tupla{nombrePromo, factor, bonoPromo}, millasGanadas:\R}$

\noindent ParametrosPromo ES $\struct{factor:\R, tope:\R}$

%ESTE ES EL NUEVO OBSERVADOR SUGERIDO - Julio
% \noindent ParametrosPromo ES $\struct{factor:\R, tope:\R, bonoPromo: \R}$

\noindent NombrePromo ES $\seq{\cha}$

\noindent Usuario ES $\Z$

\bigskip

\begin{tad}{AirMiles}
    \obs{promociones}{\dict{nombrePromo}{parametrosPromo}}
    \\
    %\obs{historial}{\dict{usuario:\Z}{\seq{Trans}}}
    \obs{historialDeTransferencias}{\seq{Trans}}
    \\
    \obs{historialDeAcumulaciones}{\dict{Usuario}{\seq{Acumulacion}}}
    \\
    \obs{historialDeCanjes}{\dict{Usuario}{\seq{millas:\Z}}}
    \\
    \obs{referidos}{\dict{Usuario}{\seq{cantidad:\Z}}}

    \bigskip
    
    %Feli
    
    \textbf{Feli}
    
    \bigskip
    
    
    \aux{usuarios}{s: AirMiles}{\conj{Usuario}}{(usuarios(s) = \conj{claves(s.historial)})}
    
    \bigskip
    
    \aux{millasDisponibles}{s: AirMiles, u: Usuario}{\R}{(\sum_{i=0}^{|s.historial[u]|-1}{s.historial[u][i].cantidad})}
    
    \bigskip
    
    %\aux{millasTotales}{s: AirMiles, u: Usuario}{\R}{(\sum_{i=0}^{|d|-1}{IfThenElse(s.historial[u].cantidad \geq 0, s.historial[u].cantidad, 0)})}
    \aux{millasTotales}{s: AirMiles, u: Usuario}{\R}{}
    
    $\bigg(\sum_{i=0}^{|s.historial[u]|-1}{IfThenElse(s.historial[u][i].cantidad \geq 0, s.historial[u][i].cantidad, 0)}\bigg)$
    
    \bigskip
    
    \aux{categoria}{s: AirMiles, u:Usuario}{Categoria}{(millasTotales(s,u) < 10.000 \implica res = "Bronce") \lor (10.000 \leq millasTotales(s,u) < 50.000 \implica res = "Plata") \lor (50.000 \leq millasTotales(s,u) < 1.000.000 \implica res = "Oro") \lor (millasTotales (s,u) \geq 1.000.000 \implies res = "Platino")}
    
    \bigskip
    
    \begin{proc}{crearSistema}{()}{\Tipo{AirMiles}}
    
        \requiere{True}
    
        \asegura{res.historial = \{\} \land res.promociones = \{"NoPromo" : (1,0)\}}
         
    \end{proc}
    
    \bigskip
    
    \begin{proc}{registrarUsuario}{\Inout s:\Tipo{AirMiles}, \In u:\Tipo{Usuario}}{}
    
        \requiere{u\notin usuarios(s) \land u \geq 0 \land s = S_0}
    
        \asegura{u \in s.historial \land setKey(s.historial, u, \struct{origen: u, destino: u, cantidad: 0, tipo: Acumular} \land (millasDisponibles(s,u) = 0 \land millasTotales(s,u) = 0) \land (\forall usu:\Z)(usu \in S_0.historial \implicaLuego (usu\in s.historial \land s.historial[usu] = S_0.historial[usu]) \land s.promociones = S_0.promociones}
    \end{proc}
    
    \bigskip
    
    
    \begin{proc}{registrarUsuarioConReferidor}{\Inout s:\Tipo{AirMiles}, \In ref:\Tipo{Usuario}, \In u:\Tipo{Usuario}}{}
    
        \requiere{u \geq 0 \land u\notin usuarios(s) \land ref \in usuarios(s) \land s=S_0}
    
        \asegura{u \in s.historial \land setKey(s.historial, u, \struct{origen: u, destino: u, cantidad: 500, tipo: Acumular} \land setKey(s.historial, ref, \struct{origen: ref, destino: ref, cantidad: 500, tipo: Acumular} \land (millasDisponibles(s,u) = 500 \land millasDisponibles(s,ref) = millasDisponibles(s_0,u) + 500) \land (millasTotales(s,u) = 500 \land millasTotales(s,ref) = millasTotales(S_0,ref) + 500) \land (\forall usu:\Z)(usu \in historial(s_0) \implicaLuego (usu\in s.historial \land s.historial[usu] = s_0.historial[usu]) \land s.promociones = S_0.promociones}
    
    \end{proc}
    
    \bigskip
    
    \begin{proc}{consultarSaldo}{\In s:\Tipo{AirMiles}, u:\Tipo{Usuario}}{\R}
    
        \requiere{u\in usuarios(s)}
    
        \asegura{res = millasDisponibles(s,u)}
         
    \end{proc}
    
    \bigskip
    
    \begin{proc}{consultarCategoria}{\In s:\Tipo{AirMiles}, u:\Tipo{Usuario}}{\Tipo{Categoria}}
    
        \requiere{u\in usuarios(s)}
    
        \asegura{res = categoria(s,u)}
    
    \end{proc}
    
    \begin{proc}{canjearMillas}{\Inout s: \Tipo{AirMiles}, u:\Tipo{Usuario}, millas:\R}{}
          
        \requiere{s = S_0 \land u\in usuarios(s) \yLuego (millas > 0 \land millas < millasDisponibles(s,u))}
          
        \asegura{s.historial[u] = s_0.historial[u] ++ \seq{\struct{origen:u, destino:u, cantidad:-millas, tipoDeTransaccion: Canjear}} \land (\forall usu:Usuario)(usu\in s_0.historial \implicaLuego (usu\in s.historial \land millasDisponibles(s, u) = millasDisponibles(S_0, u) - millas \land millasTotales(s,u) = millasTotales(s_0,u))) \land s.promociones = S_0.promociones }
          
    \end{proc}
    
    (aca quiero decir que todo usuario que estaba en el historial va a seguir estando, al igual que todas sus transacciones anteriores (solo agrego una transaccion de tipo Canjear), pero no se si esta bien escrito de esta manera y si esta bien concatenar asi para agregar una nueva transaccion, todavia me falta aprender pero por lo menos estan las ideas, se corrige y estamos. También por ahi me faltan predicados para las formulas largas,  paratodo, o que los usuarios estan o no estan en el historial (mas errores de sintaxis: como agregar una nueva transacción a la secuencia de structs, cómo seleccionar un atributo/elemento del struct?, por lo que vi en el apunte de especifcacion es .cantidad como lo hice, y para agregar a una secuencia, concateno y se agrega al final. Preguntar si esta bien. Principalmente consultas de sintaxis y correctitud. También decidir porque comparte cosas con transferirMillas, a ver si uso subSeq,soloCambianEstos y agregoTrans de tipo Canje. (MERGEAR))).
    
    \textbf{Feli}
    
    \bigskip
    %Feli
    
    \begin{proc}{rankingMillas}
    {
    \In N:\Z
    }{
    seq<\Z>
    }
    \requiere{N>0}
    \asegura{(estaOrdenadosegunCantidadDeMillasDisponibles(res,t,s) \wedge \\
         todosUsuariosEstanEnElHistorial(res,d) \wedge \\
         ordenResMayorIgualQueObs(u,s)) \wedge \\
        ((cantidadClaves.historial \geq N \wedge |res.historial|=N) \lor \\
         (cantidadClaves.historial < N \wedge |res.historial|=|cantidadClaves.historial|))}
    \end{proc}
    
    \bigskip
    
    % 1. cantidadClaves
\aux{cantidadClaves}{d: dict<\Z, seq<Trans>>}{\Z}{|d|}

% 2. estaOrdenadoSegunCantMillasDisponibles
\pred{estaOrdenadoSegunCantMillasDisponibles}{u: seq<\Z>, t: String, s: AirMiles}{
    (\forall i : \mathbb{Z}) (0 \le i < |u| - 1 \rightarrow _L sumaMillasDisponibles(u, t, s)[i] \ge sumaMillasDisponibles(u, t, s)[i+1])
}

% 3. sumaMillasDisponibles
\aux{sumaMillasDisponibles}{u: seq<\Z>, t: String, s: AirMiles}{\Z}{
    sumarTipoAcumular(u, t, s) + sumarTipoCanjear(u, t, s) + 
    sumarTipoTransferencia(u, t, s) + \\ 
    sumarTipoTransferenciaRecibida(u, t, s)
}

% 4. todosUsuariosEstanEnElHistorial
\pred{todosUsuariosEstanEnElHistorial}{l: seq<\Z>, d: dict<\Z, seq<Trans>>}{
    (\forall i : \mathbb{Z}) (0 \le i < |l| \implies l[i] \in claves(d))
}

% 5. ordenResMayorIgualQueObs
\pred{ordenResMayorIgualQueObs}{u: \seq{\Z}, s: AirMiles}{
    (\forall i: \Z)(0 \leq i < \norm{u} - 1 \implicaLuego
    \sum_{k=0}^{|claves(s)| - 1} \text{ifThenElse}(\text{SumarMillasDisponibles}(claves(s)[k], s) \geq \text{SumarMillasDisponibles}(res[i], s), 1, 0) = i + 1)
}

    \bigskip

    \newpage

    \textbf{OK}

    \begin{proc}{iniciarPromocion}
    {
        \Inout s: \Tipo{AirMiles},
        \In nombre: \seq{\cha},
        \In factor: \R
    }{}
        \requiereLargo{
            \norm{nombre} > 0 \land \\
            factor \geq 1 \land \\
            s = s_0 \land \\
            nombre \notin s.promociones
        }
        \aseguraLargo{
            seAgregaPromocion(s_0, s, nombre, \struct{factor: factor, tope: {-1.0}}) \land \\
            mismosHistorialesYReferidos(s_0, s)
        }
    \end{proc}

    \bigskip

    \textbf{OK}

    \begin{proc}{iniciarPromocionConTope}
    {
        \Inout s: \Tipo{AirMiles},
        \In nombre: \seq{\cha},
        \In factor: \R,
        \In tope: \R
    }{}
        \requiereLargo{
            \norm{nombre} > 0 \land \\
            factor \geq 1 \land \\
            tope > 0 \land \\
            s = s_0 \land \\
            nombre \notin s.promociones
        }
        \aseguraLargo{
            seAgregaPromocion(s_0, s, nombre, \struct{factor: factor, tope: tope}) \land \\
            mismosHistorialesYReferidos(s_0, s)
        }
    \end{proc}

    \bigskip

    \textbf{OK}

    \begin{proc}{finalizarPromocion}
    {
        \Inout s: \Tipo{AirMiles},
        \In nombre: \seq{\cha}
    }{}
        \requiereLargo{
            \norm{nombre} > 0 \land \\
            nombre \neq \text{``NoPromo''} \land \\
            nombre \in s.promociones \land \\
            s = s_0
        }
        \aseguraLargo{
            seEliminaPromocion(s_0, s, nombre) \land \\
            mismosHistorialesYReferidos(s_0, s)
        }
    \end{proc}

    \bigskip

    \textbf{OK}

    \predLargo{mismosHistorialesYReferidos}{s_1: \Tipo{AirMiles}, s_2: \Tipo{AirMiles}}{
        s_1.historialDeTransferencias = s_2.historialDeTransferencias \land \\
        s_1.historialDeAcumulaciones = s_2.historialDeAcumulaciones \land \\
        s_1.historialDeCanjes = s_2.historialDeCanjes \land \\
        s_1.referidos = s_2.referidos
    }

    \bigskip

    \textbf{OK}

    \predLargo{seAgregaPromocion}{s_0: \Tipo{AirMiles}, s: \Tipo{AirMiles}, nombre: \seq{\cha}, datosPromo: \Tipo{ParametrosPromo}}{
        \norm{s.promociones} = \norm{s_0.promociones} + 1 \land \\
        nombre \in s.promociones \land \\
        s.promociones[nombre] = datosPromo \land \\
        (\forall r: \seq{\cha})(r \in s_0.promociones \implicaLuego \\
        \hspace*{2em} (r \in s.promociones \yLuego s.promociones[r] = s_0.promociones[r]))
    }

    \bigskip

    \textbf{OK}

    \predLargo{seEliminaPromocion}{s_0: \Tipo{AirMiles}, s: \Tipo{AirMiles}, nombre: \seq{\cha}}{
        \norm{s_0.promociones} = \norm{s.promociones} + 1 \land \\
        nombre \notin s.promociones \land \\
        (\forall r: \seq{\cha})(r \in s.promociones \implicaLuego \\
        \hspace*{2em} (r \in s_0.promociones \yLuego s_0.promociones[r] = s.promociones[r]))
    }

    \bigskip

%-------------EMPIEZA transferirMillas JUNTO CON SUS AUX Y PREDS
    
    % transferirMillas listo
    \begin{proc}{transferirMillas}
    {
    \In origen: \Z ,
    \In destino: \Z ,
    \In monto: \R ,
    \Inout s: \Tipo{AirMiles}
    }{
    }
        \requiereLargo{
            (s = s_0) \land (origen \in s) \land (destino \in s) \land (origen \neq destino) \land
            disponeDeMillas(origen, monto, s_0[origen])
            \land (observadorValido(s))
        }
        \aseguraLargo{
        (\norm{s} = \norm{s_0}) \land (s.promociones = s_0.promociones) \land
        soloCambianEstos(\seq{origen, destino}, s, s_0) \land
        tamanoDeHistorialesCorrecto(origen, destino, s, s_0) \land
        esSubSeq(s_0[origen], s[origen]) \land
        esSubseq(s_0[destino], s[destino]) \land
        seAgregoTrans(
            \struct{origen: origen, destino: destino,
            monto: -monto, \text{``transferencia''}},
            s_0[origen], s[origen]) \land
        seAgregoTrans(
            \struct{origen: origen, destino: destino,
            monto: monto, \text{``transferencia''}},
            s_0[destino], s[destino])
        }
    \end{proc}
    
    % disponeDeMillas listo: revisa si el usuario tiene suficientes millas
    % disponibles para transferir a otro
    \predLargo{disponeDeMillas}{
    \In origen: \Z ,
    \In monto: \Z ,
    \In transDeUsuario: \seq{\struct{\Z, \Z, \R, \seq{\cha}}}
    }{
    (\sum_{i=0}^{\norm{transDeUsuario} -1} 
        transaccionesDeUsuario[i].cantidad) \geq monto 
    }
    
    \bigskip
    
    % soloCambianEstos listo: verifica que todas las transacciones de todos los 
    % usuarios que NO ESTEN en la secuencia "estos" sean iguales, garantizando
    % que solo las transacciones de "estos" sean diferentes
    
    \predLargo{soloCambianEstos}{
    \In estos: \seq{\Z} ,
    \In S: AirMiles ,
    \In S_0: AirMiles 
    }{
    (\forall usuario \in claves(S))( usuario \notin estos) \implies S[usuario] = S_0[usuario])
    }
    
%-------------TERMINA transferirMillas CON SUS AUX Y PREDS

    \begin{proc}{elQueTieneMasMillas}
    {
    \In s: \Tipo{AirMiles}
    }{
    \Z
    }
    \requiere{|claves(s.historialDeAcumulaciones)| > 0}
    \asegura{
        res \in claves(s.historialDeAcumulaciones)
        \land
        (\forall u:\mathbb{Z})
        (u \in claves(s.historialDeAcumulaciones) \implicaLuego
        millasTotales(s,u) \leq millasTotales(s,res))
    }
    \end{proc}

\aux{millasTotales}{u:Usuario, s: AirMiles}{\Z}{
    sumarTipoAcumular(u,t,s) + sumarTipoTransferenciaRecibida(u,t,s)
}

    \bigskip
 
    \begin{proc}{elMasCanjeador}
    {
    \In d: dict<\Z,Trans>,
    \In t: \Tipo{TipoDeTransaccion},
    \In s: \Tipo{AirMiles}
    }{
    \Z
    }
    \requiere{(\forall k:\mathbb{Z}) sumaMillasCanjeadas(k,t,s) \leq \\ 
    sumaMillasDisponibles(k,t,s)}
    \asegura{ res \in claves(d.historial) \wedge (\forall k:\mathbb{Z})(k \in claves(d.historial) \rightarrow _L \\ 
    sumaMillasPerdidas(k,t,s) \leq sumaMillasPerdidas(res,t,s))}
    \end{proc}
    
    \bigskip
    % --- Procedimiento Principal ---
\begin{align*}
&\textbf{proc } \text{paresTransferidoresExclusivos}(t: \text{tuple}, r: \text{string}, s: \text{Airmiles}) : \text{seq}\langle\mathbb{Z}, \mathbb{Z}\rangle \{ \\
&\quad \textbf{requiere } \{ \text{True} \} \\
&\quad \textbf{asegura } \{ \text{todosUsuariosDistintos}(res) \land |res|=|s|/2 \wedge \text{sumarTipoTransferencia}(t_0, r, s) > 0 \\
&\qquad \land \text{sumarTipoTransferencia}(t_1, r, s) > 0 \land \text{todosParesUsuariosSeTransSoloEntreSi}(t, res) \\
&\qquad \land \langle t_0, t_1 \rangle \in res \iff \langle t_1, t_0 \rangle \notin res \} \\
&\}
\end{align*}
%falta poner en el asegura que el tamaño de res es <= que la mitad de toda la lista de usuarios
% --- Predicado todosUsuariosDistintos ---
\begin{align*}
\text{pred } \text{todosUsuariosDistintos}(res: \text{seq}\langle\langle\mathbb{Z}, \mathbb{Z}\rangle\rangle) \{ \\
\quad (\forall i: \mathbb{Z})(0 \le i < |res| \rightarrow_L ((\forall j: \mathbb{Z})(0 \le j < |res| \land i \neq j \rightarrow_L \\
\quad ((res[i]_0 \neq res[j]_0 \land res[i]_1 \neq res[j]_1 \land res[i]_0 \neq res[j]_1))))) \\
\}
\end{align*}

% --- Predicado sumarTipoTransferencia ---
\begin{align*}
&\textbf{pred } \text{sumarTipoTransferencia}(u: \mathbb{Z}, r: \text{string}, s: \text{Airmiles}) \quad \% \text{ HECHO POR JULIO}
\end{align*}

% --- Predicado todosParesUsuariosSeTransSoloEntreSi ---
\begin{align*}
&\textbf{pred } \text{todosParesUsuariosSeTransSoloEntreSi}(t: \langle\mathbb{Z}, \mathbb{Z}\rangle, res: \text{seq}\langle\mathbb{Z}, \mathbb{Z}\rangle) \{ \\
&\quad t \in res \iff \text{hayTransferencia}(t_0, t_1) \land \text{hayTransferencia}(t_1, t_0) \land (\forall s: \langle\mathbb{Z}, \mathbb{Z}\rangle) \\
&\quad (s \neq t \land s \in res \to \neg \text{hayTransferencia}(t_0, s_0) \land \neg \text{hayTransferencia}(t_0, s_1) \land \\
&\quad \neg \text{hayTransferencia}(t_1, s_0) \land \neg \text{hayTransferencia}(t_1, s_1)) \\
&\}
\end{align*}

% --- Predicado hayTransferencia ---
\begin{align*}
&\textbf{pred } \text{hayTransferencia}(x: \mathbb{Z}, y: \mathbb{Z}) \{ \quad // \text{ COMPLETAR} \quad \}
\end{align*}

% los auxiliares de el mas canjeador
\aux{sumaMillasCanjeadas}{u: \Z, t: TipoDeTransaccion, s: AirMiles}{\Z}{
    sumarTipoCanjear(u,t,s)
}

\aux{sumaMillasPerdidas}{u: \Z, t: TipoDeTransaccion, s: AirMiles}{\Z}{
    sumarTipoCanjear(u,t,s) + sumarTipoTransferencia(u,t,s)
}

% Empiezan los procs de julio junto con los preds y aux especificos de
% sus procs -Julio 
\

\begin{proc}{acumularMillas}
    {
    \In u: Usuario ,
    \In millasVoladas: \R ,
    \In promo: \seq{\cha} ,
    \Inout s: AirMiles
    }{}
    {
        \requiereLargo{
            (s = s_0) \yLuego (existeUsuario(u, S)) \land
            (existePromo(promo, S)) \land
            (millasVoladas > 0)
        }
        
\aseguraLargo{
        % el unico observador que cambia es el historialDeAcumulaciones y
        % el de promociones (de forma condicional)
        (noCambiaElHistorialDeCanjes(s_0,s))
        \land (noCambiaElHistorialDeReferidos(s_0,s))
        \land (noCambiaElHistorialTransferencias(s_0,s))
        \land esSubSeq(s_0.historialDeAcumulaciones[u], s.historialDeAcumulaciones[u])
        % la siguiente expresion tiene esta Estructura
        % ((P(a,b) land c) lor (no P(a,b) land d)) la cual es
        % exactamente lo que hace un ifThenElse pero usando solo logica.
        % si se cumple P(a,b), hace c
        % si se cumple no P(a,b) (el caso contrario a P(a,b) que seria el else), hace d
        % lo que equivale a if(P(a,b){ c } else { d })
        % hice esto para no especificar COMO HACERLO (con ifThenElse)
        % y especificar QUE HACE
        % la siguiente expresion expresion es 1 solo predicado, o sea Q(x)
        % empieza el if
        \land
        (
            (
                (millasVoladas * s_0.promociones[promo].factor)
                \geq (s_0.promociones[promo].tope)
                \land (\norm{s.promociones} = \norm{s_0.promociones} - 1)
                \land s.promociones = listaDePromocionSinP(promo, s_0)
                \land seAgregoVuelo(crearAcumulacion(u, millasVoladas, promo,
                    s_0.promociones[promo].factor,
                    s_0.promociones[promo].tope,
                    s_0.promociones[promo].tope + (millasVoladas * factorDeCategoriaUsuario(u, s_0)))
                )
            )
            \lor
            (
                % empieza el else
                (millasVoladas * s_0.promociones[promo].factor)
                < (s_0.promociones[promo].tope)
                \land (\norm{s.promociones} = \norm{s_0.promociones})
                \land (soloCambioLaPromoP(promo, s, s_0))
                \land s.promociones[promo] = reducirTopeAPromoP(promo, s_0)
                \land \norm{s.historialDeAcumulaciones[u]}
                        = \norm{s_0.historialDeAcumulaciones[u]} + 1
                \land seAgregoVuelo(crearAcumulacion(u, millasVoladas, promo,
                    s_0.promociones[promo].factor,
                    millasVoladas * s_0.promociones[promo].factor,
                    (millasVoladas * factorDeCategoriaUsuario(u, s_0)
                        * s_0.promociones[promo].factor))
                )
            )
        )
        }
        
    }
\end{proc}

% EL PROC acumularMillas ESTA INCOMPLETO, FALTA EL CASO CUANDO EL USUARIO
% NO UTILIZA PROMO (utiliza la promo     noPromo) Y HAY QUE AGREGAR UNA
% EXPRESION ANALOGA HORRIBLE COMO LA FINAL DEL PROC acumularMillas
% A MENOS QUE SE UTILICE ifThenElse

\

\pred{cantidadDeVuelosCorrecto}{
\In u: Usuario, 
\In s_0: AirMiles, 
\In s: AirMiles}{
    \norm{s[u]} = \norm{s_0[u]} + 1
}

\

% pred que quiza es util para alejandro

\auxLargo{listaDePromocionSinP}{
\In promo: \seq{\cha},
\In s: AirMiles
}{
\dict{\seq{\cha}}{ParametrosPromo}
}{
    \paraTodo{p}{claves(s.promociones)}{(promo \neq p) \implies 
        (res.p = s.promociones[p])}    
}

\

% pred que quiza es util para alejandro

\predLargo{soloCambioLaPromoP}{
\In promo: \seq{char},
\In s: AirMiles,
\In s_0: AirMiles
}{
    (s.promociones[promo] \neq s_0.promociones[promo]) \\
    \land \paraTodo{p}{claves(s_0.promociones)}{(promo \neq p) \implies
        (s.promociones[p] = s_0.promociones[p])}
}

\

% pred que quiza es util para felipe

\predLargo{seAgregoVuelo}{
\In vuelo: Acumulaciones,
\In s_0: \seq{Acumulaciones},
\In s: \seq{Acumulaciones}
}{
    % en este caso no utilizo el yLuego para que no se indefina, lo utilizo
    % para que se evalue primero que el vuelo existe en s y despues que se
    % que existe, comparo las apariciones de vuelos en las dos secuencias
    % para que sea la correcta, verificar esto con el profe si
    % (vuelo \in s) deberia de ir en el asegura del proc acumularMillas
    % en lugar de un pred <- Preguntar al profesor -Julio
    (vuelo \in s) \yLuego
    (aparicionesEnSeqDe(vuelo, s) = aparicionesEnSeqDe(vuelo, s_0) + 1)
}

\

% pred que quiza es util para felipe

\auxLargo{crearAcumulacion}{
\In u: Usuario,
\In distancia: \R,
\In nombrePromo: NombrePromo,
\In factor: \R,
\In millasPromo: \R,
\In millasGanadas: \R
}{Acumulacion}{
    % crea el struct por extension
    \struct{
        "distanciaVolada": distancia, \\
        "promo": \tupla{nombrePromo, factor, millasPromo} \\
        "millasGanadas": millasGanadas
    }
}

\

% pred definitivamente util para consultarCategoria

\auxLargo{factorDeCategoriaUsuario}{\In u: Usuario, \In s: AirMiles}{\R}{
    1 + factorPlata(u,s) + factorOro(u,s) + factorPlatino(u,s)
}

\

\auxLargo{factorPlata}{
\In u: Usuario,
\In s: AirMiles
}{\R}{
    ifThenElse(sumarMillasTotalesDeUsuario(u,s) \geq 10000 , 0.5, 0)
}

\

\auxLargo{factorOro}{
\In u: Usuario,
\In s: AirMiles
}{\R}{
    ifThenElse(sumarMillasTotalesDeUsuario(u,s) \geq 50000 , 0.5, 0)
}
\

\auxLargo{factorPlatino}{
\In u: Usuario,
\In s: AirMiles
}{\R}{
    ifThenElse(sumarMillasTotalesDeUsuario(u,s) \geq 1000000 , -2, 0)
}

\

\begin{proc}{resumenDeTransacciones}{
\In u: Usuario,
\In s: AirMiles
}{\struct{total: \R, canjeadas: \R, enviadas: \R, recibidas: \R}}
{
    \requiere{existeUsuario(u,s)}
    \aseguraLargo{
        (res.total = sumarMillasTotalesDeUsuario(u,s)) \land
        (res.canjeadas = sumarCanjeadoPorUsuario(u,s)) \land
        (res.enviadas = sumarEnviadoPorUsuario(u,s)) \land
        (res.recibidas = sumarRecibidoPorUsuario(u,s))
    }
}

\end{proc}

\

\begin{proc}{transferirMillas}
    {
    \In origen: \Z ,
    \In destino: \Z ,
    \In monto: \R ,
    \Inout S: \Tipo{AirMiles}
    }{
    }
        \requiereLargo{
            (s = s_0) \yLuego (existeUsuario(origen, s_0)) \land
            (existeUsuario(destino, s_0)) \land
            (origen \neq destino) \land
            disponeDeMillas(origen, monto, s_0)
        }
        \aseguraLargo{
        % el unico observador que cambia es el historialDeTransferencias
        (noCambiaElHistorialDePromos(s_0,s)) \land
        (noCambiaElHistorialDeAcumulaciones(s_0,s)) \land
        (noCambiaElHistorialDeCanjes(s_0,s)) \land
        (noCambiaElHistorialDeReferidos(s_0,s)) \land
        tamanoDeHistorialDeTransCorrecto(s, s_0) \land
        esSubSeq(s_0.historialDeTransferencias, s.historialDeTransferencias) \land
        seAgregoTrans(
            \struct{origen: origen, destino: destino, cantidad: monto},
            s_0.historialDeTransferencias,
            s.historialDeTransferencias
        )
        }
    \end{proc}

\

\predLargo{tamanoDeHistorialDeTransCorrecto}{\In s: Airmiles, \In s_0: Airmiles}{
    \norm{S.historialDeTransferencias} = \norm{S_0.historialDeTransferencias} + 1
}

\

\predLargo{seAgregoTrans}{
\In trans: Trans,
\In S_0: historialDeTransferencias,
\In S : historialDeTransferencias
}{
    % hace falta un operador Yluego que verifique primero que la trans existe
    % en S para luego calcular cuantas veces ocurre? o los operadores 
    % LUEGO solo se utilizan exclusivamente para que no se indefina una 
    % expresion logica? <- para preguntar en la practica - Julio
    (trans \in S) \\
    \land ((aparicionesEnSeqDe(trans, S_0) + 1 = aparicionesEnSeqDe(trans, S)))
}

\

% Terminan los procs de julio junto con los preds y aux especificos de
% sus procs -Julio 


%  EMPIEZAN NUEVOS AUX Y PREDS GENERALES UTILES PARA LOS NUEVOS OBSERVADORES 
%  QUE DEFINIMOS EL VIERNES 4/9/2026 EN EL LABO -Julio

\

% quiza verificar que |S0| <= |S| es sobre especificar, pues si el proc
% que utiliza esSubSeq tambien verifica en el requiere algo como |S0| <= |S|
% se repite 2 veces la expresion |S0| <= |S|, para pensar -Julio

\predLargo{esSubSeq}{
\In S_0: \seq{T},
\In S: \seq{T}
}{
    (\norm{S_0} \leq \norm{S}) \\
    \yLuego \paraTodo{i}{\Z}{ (0 \leq i < \norm{S_0}) \implicaLuego (S_0[i] = S[i])}
}

\

% util para el requiere de muchos procs

\pred{existeUsuario}{
\In u: Usuario,
\In s: AirMiles
}{
    u \in s.historialDeAcumulaciones
}

\

% util para el requiere de muchos procs

\pred{existePromo}{
\In promo: NombrePromo,
\In s: AirMiles
}{
    (promo \in s.promociones)
}

\

% util para el asegura de muchos procs que manipula el sistema

\predLargo{noCambiaElHistorialDePromos}{
\In s_0: AirMiles,
\In s: AirMiles
}{
    (s_0.promociones = s.promociones)
}

\

% util para el asegura de muchos procs que manipula el sistema

\predLargo{noCambiaElHistorialTransferencias}{
\In s_0: AirMiles,
\In s: AirMiles
}{
    (s_0.historialDeTransferencias = s.historialDeTransferencias)
}

\

% util para el asegura de muchos procs que manipula el sistema

\predLargo{noCambiaElHistorialDeAcumulaciones}{
\In s_0: AirMiles,
\In s: AirMiles
}{
    (s_0.historialDeAcumulaciones = s.historialDeAcumulaciones)
}

\

% util para el asegura de muchos procs que manipula el sistema

\predLargo{noCambiaElHistorialDeCanjes}{
\In s_0: AirMiles,
\In s: AirMiles
}{
    (s_0.historialDeCanjes = s.historialDeCanjes)
}

\

% util para el asegura de muchos procs que manipula el sistema

\predLargo{noCambiaElHistorialDeReferidos}{
\In s_0: AirMiles,
\In s: AirMiles
}{
    (s_0.referidos = s.referidos)
}

\

% Util para los procs transferirMillas y canjearMillas

\predLargo{disponeDeMillas}{
\In u: Usuario,
\In millas: \R,
\In s: AirMiles
}{
    sumarMillasTotalesDeUsuario(u,s) \geq millas
}

\

% verifica cuantas veces se repite el elemento e de tipo T (lo que sea)
% en una secuencia de tipo T (mismo tipo que e)
\auxLargo{aparicionesEnSeqDe}{
\In e: T,
\In s: \seq{T}
}{
\Z
}{
\sum_{i=0}^{\norm{s} -1}{ifThenElse(s[i] = e , 1, 0)}
}

\

% Util en general

\auxLargo{sumarMillasTotalesDeUsuario}{
\In u: Usuario,
\In s: AirMiles
}{
\R
}{
    \sum_{i = 0}^{\norm{s.historialDeAcumulaciones[u]} -1}{
    s.historialDeAcumulaciones[u][i].millasGanadas
    } \\
    + \\
    \sum_{i = 0}^{\norm{s.referidos[u]} -1}{
    s.referidos[u][i].cantidad
    }
    % la suma de los referidos no lo piden en ningun proc, asi que no vale
    % la pena hacerlo su propio aux
}

\

% util para elMasCanjeador y resumenDeTransacciones

\auxLargo{sumarCanjeadoPorUsuario}{
\In u: Usuario,
\In s: AirMiles
}{
\R
}{
    \sum_{i = 0}^{\norm{s.historialDeCanjes[u]} -1}{
    s.historialDeCanjes[u][i]
    }
}

\

% quiza es especifico de un proc, hacer chequeo luego

\auxLargo{sumarEnviadoPorUsuario}{
\In u: Usuario,
\In s: AirMiles
}{
\R
}{
    % hay que echarle un ojo a esto para ver si queda mas prolijo en el pdf
    % - Julio
    \sum_{i = 0}^{\norm{s.historialDeTransferencias[u]} -1}{} \\
    ifThenElse(
        s.historialDeTransferencias[i].origen = u , \\
        s.historialDeTransferencias[i].cantidad , \\
        0)
}

\

% quiza es especifico de un proc, hacer chequeo luego

\auxLargo{sumarRecibidoPorUsuario}{
\In u: Usuario,
\In s: AirMiles
}{
\R
}{
    % hay que echarle un ojo a esto para ver si queda mas prolijo en el pdf
    % - Julio
    \sum_{i = 0}^{\norm{s.historialDeTransferencias[u]} -1}{} \\
    ifThenElse(
        s.historialDeTransferencias[i].destino = u , \\
        s.historialDeTransferencias[i].cantidad , \\
        0)
}

% util en general para el proc del ranking, elQueTieneMasMillas y quiza otros mas

\auxLargo{sumarMillasTotalesDeUsuario}{
\In u: Usuario,
\In s: AirMiles
}{
\R
}{
    sumarMillasTotalesDeUsuario(u,s) \\
    + sumarRecibidoPorUsuario(u,s) \\
    - sumarEnviadoPorUsuario(u,s) \\
    - sumarCanjeadoPorUsuario(u,s)
}

%  TERMINAN NUEVOS AUX Y PREDS GENERALES UTILES PARA LOS NUEVOS OBSERVADORES 
%  QUE DEFINIMOS EL VIERNES 4/9/2026 EN EL LABO -Julio

\end{tad} 

\section{borrador}

Suposiciones:

Si acumulas o canjeas, el destino y origen es uno mismo

tope -1 es decir que no tiene tope

Suponemos que ganar por referencia = ganar por viaje -> transaccionAcumular.

millas acumuladas/totales: todo lo que ganaste
millas disponibles/actuales: son las actuales

\begin{itemize}
\item tipoDeTransaccion = enum(transferencia, acumular, canjear)

\item categorias = enum(bronce,plata,oro,platino)
\end{itemize}


Observadores:

promociones dicc<nombre: seq<char>. Tupla(factor:Z, tope:int)>

historialTransacciones $dict<clave=idUsuario:Z , valor=transacciones:seq<struct<origen:Z,destino:Z,cantidad:R,tipoDeTransaccion:seq<char>>>$


Ejemplo/Testing:

constructor crearSistema()

registrarUsuario(ID) -> historialTransaccion(id, 0,0,0,0)

registrarUsuarioConReferidor -> historialTransaccion(id, 500,...)

\end{document}
